\section{Related work}
\label{sec:related-work}

Formalizing variable binding has a long history across proof
assistants, and several general-purpose approaches exist alongside
the direct, elementary style used in CSLib. Nominal
techniques~\cite{pitts2013a}, as implemented in Nominal
Isabelle~\cite{urban2005a}, provide binder-generic reasoning
principles via permutations and freshness, at the cost of adopting a
dedicated meta-theory. Locally nameless
representations~\cite{aydemir2008a} aim to make binder-heavy formalizations more
tractable or more automatically generated from a specification,
respectively. Higher-order abstract syntax sidesteps binder
representation issues by delegating them to the proof assistant's
own binding mechanism, again at the cost of restricting which proof
techniques remain available (e.g.\ induction principles).

CSLib's current approach is deliberately more elementary: named terms
and de Bruijn terms are both represented directly as first-order
inductive types, with the relationships between representations
(such as the isomorphism of Section~\ref{sec:de-bruijn}) proved
explicitly rather than assumed via a shared meta-level binding
mechanism. This trades some proof automation for representations
that are simple to state, and results that transfer cleanly between
representations by construction, which fits CSLib's goal of a
reusable foundation rather than a self-contained framework.

